% =========================================================================
% Capítulo 3: Cinética de sistemas de partículas
% Tema 01: Sistemas de partículas
% Problema 06: Mancuerna en bloque con pista circular
% Origen: Ayudantía 5 (Sólido rígido), 2026
% =========================================================================
\subsubsection{Mancuerna en bloque con pista circular}

Una mancuerna constituida por dos partículas $A$ y $B$, de masa $m$ cada una, se desliza sin fricción a lo largo de una pista circular de radio $R$ cortada en el interior de un bloque de masa $M$. El bloque, a su vez, se desliza sin fricción a lo largo de una superficie horizontal. Sabiendo que $\theta$ es el ángulo formado por la vertical y la línea que une el centro de la pista $O$ con el centro de la barra $C$, y que $x$ representa el desplazamiento horizontal del bloque, realice lo siguiente:

\begin{enumerate}
    \item Defina los marcos de referencia y sistemas de coordenadas adecuados (fijo y móvil).
    \item Analice la cinemática del sistema para expresar la posición y velocidad de las partículas en función de las coordenadas generalizadas $x$ y $\theta$.
    \item Analice la cinética del sistema para determinar las ecuaciones diferenciales del movimiento utilizando las leyes de Newton (balance de momentum lineal y angular).
\end{enumerate}

\begin{figure}[!ht]
    \centering
    \begin{tikzpicture}[>=Latex, scale=1.4]
        % --- Parámetros ---
        \def\R{2.2}       % Radio de la pista
        \def\halfL{1.4}   % Mitad de la longitud de la barra (l)
        \def\distA{1.7}   % Distancia a (ajustada visualmente < R)
        \def\ang{35}      % Ángulo theta
        \def\BlockSize{3.0} % Tamaño medio del bloque

        % Coordenada Origen del bloque (centro del círculo)
        \coordinate (O) at (0,0);

        % --- 1. Suelo ---
        \draw[thick] (-4, -\BlockSize) -- (4, -\BlockSize);
        \fill[pattern=north east lines] (-4, -\BlockSize-0.2) rectangle (4, -\BlockSize);

        % --- 2. Bloque M ---
        % El bloque es un rectángulo con un agujero circular
        \fill[gray!5, even odd rule] 
            (-\BlockSize, -\BlockSize) rectangle (\BlockSize, \BlockSize)
            (0,0) circle (\R);
        \draw[thick] (-\BlockSize, -\BlockSize) rectangle (\BlockSize, \BlockSize);
        \draw[thick] (0,0) circle (\R);
        
        % Hachurado en el borde interior de la pista (estilo técnico)
        \path[pattern=north east lines, pattern color=gray!80, even odd rule] (0,0) circle (\R+0.2) circle (\R+0.4);
        \filldraw[white, draw= black] (0,0) circle (\R+0.2);
        
        % Etiqueta del Bloque
        \node at (-\BlockSize+0.5, \BlockSize-0.5) {$M$};

        % --- 3. Elementos Geométricos (Internos) ---
        \begin{scope}[rotate=-\ang] 
            % Línea OC (distancia a)
            \draw[dashed] (O) -- (0, -\distA) node[midway, left] {$a$};
            \coordinate (C) at (0, -\distA);
            \filldraw (C) circle (1pt) node[below left, font=\footnotesize] {$C$};

            % Barra AB (Perpendicular a OC)
            \draw[line width=3pt, gray!40] (-\halfL, -\distA) -- (\halfL, -\distA);
            \draw[thick] (-\halfL, -\distA) -- (\halfL, -\distA);

            % Partículas
            \coordinate (A) at (-\halfL, -\distA);
            \coordinate (B) at (\halfL, -\distA);

            \shadedraw[ball color=gray!40] (A) circle (0.2) node[above right, xshift=5pt] {$A$};
            \shadedraw[ball color=gray!40] (B) circle (0.2) node[above right, xshift=5pt] {$B$};
            \node[ font=\footnotesize] at (A) {$m$};
            \node[ font=\footnotesize] at (B) {$m$};

            % Radio R (hacia B, auxiliar para mostrar geometría)
            \draw[->, thin, gray] (O) -- (0,\R+0.2) node[midway,below right, black] {$R$};
            
            % Cota 2l
            \draw[|<->|] ($(A)+(0,-0.3)$) -- ($(B)+(0,-0.3)$) node[midway, below] {$2l$};
        \end{scope}

        % --- 4. Referencias y Cotas Externas ---
        % Centro O
        \filldraw (O) circle (1.5pt) node[above left] {$O$};
        
        % Línea vertical referencia para theta
        \draw[dashed, gray] (O) -- (0, -2.0);
        \draw[->] (0, -0.8) arc (-90:{-90-\ang}:0.8);
        \node at (-0.2, -1.1) {$\theta$};

        % Coordenada x del bloque
        \draw[|->, thick] (0.5, -\BlockSize-0.5) -- (1.5, -\BlockSize-0.5) node[midway, above] {$x$};
        
        % Gravedad
        \draw[->, thick] (3.5, 0.5) -- (3.5, -0.5) node[midway, right] {$\mathbf{g}$};

    \end{tikzpicture}
    \caption{Sistema de mancuerna deslizante dentro de un bloque móvil con cavidad circular.}
    \label{fig:mancuerna_bloque_movil}
\end{figure}

% --- SOLUCIÓN ---------------------------------------------------------------
\begin{solucion}{Mancuerna en bloque con pista circular (por Newton)}
\textbf{Marcos de referencia y Sistemas de Coordenadas}\\
Para este caso podemos darnos cuenta de que tenemos dos cuerpos rígidos, el bloque y la mancuerna.
\begin{itemize}
    \item $\mathcal{N}$: suelo (marco inercial) $\longrightarrow\{\vec{e}_x;\vec{e}_y;\vec{e}_z\}$ cartesiano fijo en $\mathcal{N}$.
    \item $\mathcal{A}$: barra / mancuerna $\longrightarrow\{\vec{e}_r;\vec{e}_\theta;\vec{e}_z\}$ cilíndrico, fijo en la barra y usado para $A$ y $B$.
\end{itemize}

\begin{center}
\begin{tikzpicture}[>=Latex, scale=1.2]
    % --- PARÁMETROS ---
    \def\R{1.6}    % radio de la pista [u]
    \def\ang{35}   % ángulo theta [°]
    \def\len{1.0}  % largo de los vectores [u]

    % --- PISTA Y CENTRO ---
    \draw[thick] (0,0) circle (\R);
    \filldraw (0,0) circle (1.2pt) node[above left] {$O$};
    \draw[dashed, gray] (0,0) -- (0,-\R);

    % --- BASE FIJA (e_x hacia abajo, e_y hacia la derecha) ---
    \draw[->, thick, myred] (0,0) -- (0,-\len) node[below left] {$\vec{e}_x$};
    \draw[->, thick, myred] (0,0) -- (\len,0) node[right] {$\vec{e}_y$};

    % --- BASE MÓVIL ---
    \draw[->, thick, azul] (0,0) -- ({-90+\ang}:\len) node[below right] {$\vec{e}_r$};
    \draw[->, thick, azul] (0,0) -- ({\ang}:\len) node[above right] {$\vec{e}_\theta$};
    \draw[dashed] (0,0) -- ({-90+\ang}:\R);

    % --- ÁNGULO ---
    \draw[->] (0,-0.5) arc (-90:{-90+\ang}:0.5);
    \node at ({-90+\ang/2}:0.72) {$\theta$};
\end{tikzpicture}
\end{center}

\subsubsection*{Análisis cinemático}
\textbf{Bloque}
$$\marc{\vec{r}_O}{N}{} = x\,\vec{e}_y, \qquad
\marc{\vec{v}_O}{N}{} = \dot{x}\,\vec{e}_y, \qquad
\marc{\vec{a}_O}{N}{} = \ddot{x}\,\vec{e}_y$$

\textbf{Partícula $A$}
$$\begin{aligned}
\marc{\vec{r}_A}{N}{} &= \marc{\vec{r}_O}{N}{}+\vec{r}_{A/O} = x\,\vec{e}_y + a\,\vec{e}_r - \ell\,\vec{e}_\theta,\\
\marc{\vec{v}_A}{N}{} &= \dot{x}\,\vec{e}_y + a\dot{\theta}\,\vec{e}_\theta - \ell\dot{\theta}\,\vec{e}_r,\\
\marc{\vec{a}_A}{N}{} &= \ddot{x}\,\vec{e}_y + \left(-a\dot{\theta}^2+\ell\ddot{\theta}\right)\vec{e}_r + \left(a\ddot{\theta}+\ell\dot{\theta}^2\right)\vec{e}_\theta.
\end{aligned}$$

\textbf{Partícula $B$}
$$\begin{aligned}
\marc{\vec{r}_B}{N}{} &= x\,\vec{e}_y + a\,\vec{e}_r + \ell\,\vec{e}_\theta,\\
\marc{\vec{v}_B}{N}{} &= \dot{x}\,\vec{e}_y + a\dot{\theta}\,\vec{e}_\theta + \ell\dot{\theta}\,\vec{e}_r,\\
\marc{\vec{a}_B}{N}{} &= \ddot{x}\,\vec{e}_y + \left(-a\dot{\theta}^2-\ell\ddot{\theta}\right)\vec{e}_r + \left(a\ddot{\theta}-\ell\dot{\theta}^2\right)\vec{e}_\theta.
\end{aligned}$$

\textbf{Otra forma de obtener las mismas aceleraciones: matrices de rotación}\\
En el espacio tridimensional, la orientación de un sistema de referencia respecto a un marco inercial fijo $\mathcal{N}$ se describe mediante una matriz de rotación ortogonal $\mathbf{R} \in SO(3)$. Una de sus propiedades fundamentales es que su derivada temporal multiplicada por su transpuesta define la matriz antisimétrica de la velocidad angular:
$$\dot{\mathbf{R}}\mathbf{R}^T = [\boldsymbol{\omega}]_\times, \qquad
[\boldsymbol{\omega}]_\times = \begin{bmatrix} 0 & -\omega_z & \omega_y \\ \omega_z & 0 & -\omega_x \\ -\omega_y & \omega_x & 0 \end{bmatrix}$$
Multiplicando por la derecha por $\mathbf{R}$ y sabiendo que $\mathbf{R}^T\mathbf{R} = \mathbf{I}$, obtenemos $\dot{\mathbf{R}} = [\boldsymbol{\omega}]_\times \mathbf{R}$. Para la segunda derivada diferenciamos con la regla del producto, y con $\boldsymbol{\alpha} = \dot{\boldsymbol{\omega}}$:
$$\ddot{\mathbf{R}} = \left( [\boldsymbol{\alpha}]_\times + [\boldsymbol{\omega}]_\times^2 \right) \mathbf{R}$$
Para una rotación exclusivamente alrededor del eje inercial $z$ con ángulo variable $\theta(t)$, $\boldsymbol{\omega} = [0, 0, \dot{\theta}]^T$ y $\boldsymbol{\alpha} = [0, 0, \ddot{\theta}]^T$, de modo que
$$[\boldsymbol{\omega}]_\times = \begin{bmatrix} 0 & -\dot{\theta} & 0 \\ \dot{\theta} & 0 & 0 \\ 0 & 0 & 0 \end{bmatrix}, \quad
[\boldsymbol{\alpha}]_\times = \begin{bmatrix} 0 & -\ddot{\theta} & 0 \\ \ddot{\theta} & 0 & 0 \\ 0 & 0 & 0 \end{bmatrix}, \quad
[\boldsymbol{\omega}]_\times^2 = \begin{bmatrix} -\dot{\theta}^2 & 0 & 0 \\ 0 & -\dot{\theta}^2 & 0 \\ 0 & 0 & 0 \end{bmatrix}$$
Con $\mathbf{r}_O = x\,\mathbf{e}_y$ y la posición de cada partícula en el marco local del cuerpo ($\mathbf{r}'_{A/O}=[a,-\ell,0]^T$ y $\mathbf{r}'_{B/O}=[a,\ell,0]^T$), para $P=A,B$:
$$\mathbf{r}_P = \mathbf{r}_O + \mathbf{R}\,\mathbf{r}'_{P/O}, \qquad
\mathbf{v}_P = \mathbf{v}_O + [\boldsymbol{\omega}]_\times \mathbf{R}\,\mathbf{r}'_{P/O}, \qquad
\mathbf{a}_P = \mathbf{a}_O + \left( [\boldsymbol{\alpha}]_\times + [\boldsymbol{\omega}]_\times^2 \right) \mathbf{R}\,\mathbf{r}'_{P/O}$$

\subsubsection*{Análisis cinético}
\begin{center}
\begin{tikzpicture}[>=Latex, scale=0.75]
    % --- PARÁMETROS ---
    \def\R{2.2}        % radio de la pista [u]
    \def\halfL{1.4}    % mitad de la varilla [u]
    \def\distA{1.7}    % distancia a [u]
    \def\ang{35}       % ángulo theta [°]
    \def\BlockSize{3.0}% semilado del bloque [u]

    % --- DCL DEL BLOQUE M ---
    \begin{scope}
        \node[font=\bfseries] at (0, \BlockSize + 0.8) {DCL: Bloque $M$};
        \draw[thick, gray!70] (-\BlockSize, -\BlockSize) rectangle (\BlockSize, \BlockSize);
        \draw[thick, gray!70] (0,0) circle (\R);
        \filldraw (0,0) circle (1.5pt) node[above left] {$O$};
        \draw[->, gray, dashed] (0,0) -- (0, -1.5) node[right] {$\vec{e}_x$};
        \draw[->, gray, dashed] (0,0) -- (1.5, 0) node[above] {$\vec{e}_y$};
        \draw[thick, ->, myred] (0,0) -- (0, -2.2) node[right] {$\vec{F}_g^O$};
        \draw[thick, ->, myred] (-0.5, -\BlockSize) -- (-0.5, -\BlockSize+2) node[left] {$\vec{N}_O$};
        \begin{scope}[rotate=-\ang]
            \coordinate (posA) at (-\halfL, -\distA);
            \coordinate (posB) at (\halfL, -\distA);
        \end{scope}
        \draw[thick, ->, myred] (posA) -- ($(posA)!-1.3cm!(0,0)$) node[left] {$-\vec{N}_A$};
        \draw[dashed, gray] (0,0) -- (posA);
        \draw[thick, ->, myred] (posB) -- ($(posB)!-1.3cm!(0,0)$) node[right] {$-\vec{N}_B$};
        \draw[dashed, gray] (0,0) -- (posB);
    \end{scope}

    % --- DCL DE LA PARTÍCULA A ---
    \begin{scope}[xshift=-6cm, yshift=-8cm]
        \node[font=\bfseries] at (0, 2.5) {DCL: Partícula $A$};
        \coordinate (A) at (0,0);
        \shadedraw[ball color=gray!40] (A) circle (0.3) node[above right, xshift=5pt] {$A$};
        \draw[thick, ->, myred] (A) -- (0, -2) node[right] {$\vec{F}_g^A$};
        \begin{scope}[rotate=-\ang]
            \coordinate (dirOA) at (\halfL, \distA);
            \draw[dashed, gray] (A) -- ($(A)!2cm!(dirOA)$) node[above] {hacia $O$};
            \draw[thick, ->, myred] (A) -- ($(A)!1.5cm!(dirOA)$) node[right] {$\vec{N}_A$};
            \draw[dashed, gray] (1.5, 0) -- (-2, 0);
            \draw[thick, ->, myred] (A) -- (-1.8, 0) node[above] {$\vec{R}_A$};
        \end{scope}
    \end{scope}

    % --- DCL DE LA PARTÍCULA B ---
    \begin{scope}[xshift=6cm, yshift=-8cm]
        \node[font=\bfseries] at (0, 2.5) {DCL: Partícula $B$};
        \coordinate (B) at (0,0);
        \shadedraw[ball color=gray!40] (B) circle (0.3) node[above left, xshift=-5pt] {$B$};
        \draw[thick, ->, myred] (B) -- (0, -2) node[right] {$\vec{F}_g^B$};
        \begin{scope}[rotate=-\ang]
            \coordinate (dirOB) at (-\halfL, \distA);
            \draw[dashed, gray] (B) -- ($(B)!2cm!(dirOB)$) node[below right] {hacia $O$};
            \draw[thick, ->, myred] (B) -- ($(B)!1.5cm!(dirOB)$) node[left] {$\vec{N}_B$};
            \draw[dashed, gray] (1.5, 0) -- (-2, 0);
            \draw[thick, ->, myred] (B) -- (1.8, 0) node[above] {$\vec{R}_B$};
        \end{scope}
    \end{scope}
\end{tikzpicture}
\end{center}

\textbf{Conjunto de fuerzas externas}\\
En este caso, al analizar el bloque, tiene 4 fuerzas: el peso, ejercido por la tierra; la normal ejercida por la superficie; y las dos normales que ejercen las masas de la mancuerna.

\textbf{Bloque}
$$\begin{aligned}
\text{a) }&\vec{F}^{\,O}_g = M g\,\vec{e}_x, \qquad
\text{b) } \vec{N}_O = -N_O\,\vec{e}_x,\\
\text{c) }&-\vec{N}_A = N_A\,\frac{a\vec{e}_x-\ell\vec{e}_y}{R}, \qquad
\text{d) } -\vec{N}_B = N_B\,\frac{a\vec{e}_x+\ell\vec{e}_y}{R}.
\end{aligned}$$
Para conocer la dirección de la fuerza normal de la mancuerna nos podemos recordar que es un círculo, por ende tiene una dirección radial; para la partícula usamos la posición de la partícula relativa al punto $O$ normalizada.

\textbf{Partícula $A$}
$$\begin{aligned}
\text{a) }&\vec{F}^{\,A}_g = m g\,\vec{e}_x = m g\cos\theta\,\vec{e}_r - m g\sin\theta\,\vec{e}_\theta,\\
\text{b) }&\vec{N}_A = -N_A\,\frac{a\vec{e}_x-\ell\vec{e}_y}{R},\qquad
\text{c) } \vec{R}_A= R\,\vec{e}_\theta.
\end{aligned}$$
Mismo análisis para la partícula $B$:
$$\begin{aligned}
\text{a) }&\vec{F}^{\,B}_g = m g\,\vec{e}_x = m g\cos\theta\,\vec{e}_r - m g\sin\theta\,\vec{e}_\theta,\\
\text{b) }&\vec{N}_B = -N_B\,\frac{a\vec{e}_x+\ell\vec{e}_y}{R},\qquad
\text{c) } \vec{R}_B= -R\,\vec{e}_\theta.
\end{aligned}$$

Ahora podemos hacer Newton, pero tenemos muchos Newton por hacer. Podemos hacer varios DCL (diagramas de cuerpo libre); ¿cuál tiene ventaja sobre otro? Aquel en el que conozcamos todas las fuerzas sobre ese eje. En este caso el sistema completo tiene 0 fuerzas externas en la dirección $\vec{e}_y$.

\textbf{Balance de fuerzas del CM en $\vec{e}_y$ para $(A+B+\text{bloque})$}
$$0 = \left( m\,\marc{\vec{a}_A}{N}{} + m\,\marc{\vec{a}_B}{N}{} + M\,\marc{\vec{a}_O}{N}{} \right)\!\cdot\!\vec{e}_y,$$
$$0 = (2m+M)\,\ddot{x} - 2ma\,\dot{\theta}^2\sin\theta + 2ma\,\ddot{\theta}\cos\theta,$$
\respuesta{$\displaystyle \ddot{x}=\frac{2ma}{2m+M}\Big(\dot{\theta}^2\sin\theta-\ddot{\theta}\cos\theta\Big)$}

Ahora necesitamos otra ecuación. ¿Cómo la obtenemos? De otro balance de fuerzas, pero en otra dirección o sobre otro sistema. Lamentablemente, en todo otro sistema hay fuerzas que no conocemos; podemos despejarlas y así resolver, pero hay una manera más fácil que en este caso calza por ser circular: si calculamos el momento angular de la mancuerna, las únicas fuerzas son el peso y las normales, y en el balance de momento angular respecto a $O$ la dirección de $\vec{r}_{A\text{ o }B}$ es la misma que la de su normal, por ende su producto cruz es 0. Esto ya viene con muchos años de análisis, pero es bueno que comiencen a trabajar para tener ese entendimiento de cuáles caminos son más fáciles (de todos modos cualquier forma de obtenerlo es válida, solo que es medio complejo). Ahora, si lo hacemos respecto a $O$ tenemos un problema: es un punto que acelera, por ende nos sale el cochino término de la aceleración del punto.

\textbf{Ecuación de torques respecto de $O$ para el sistema $(A+B)$}
$$\der{\marc{\vec{H}_O}{N}{}}{N}{} = 2mR^2\,\ddot{\theta}\,\vec{e}_z,$$
$$\begin{aligned}
\marc{\vec{M}_O}{N}{}
&= (\vec{r}_A-\vec{r}_O)\times\vec{F}^{\,A}_g + (\vec{r}_B-\vec{r}_O)\times\vec{F}^{\,B}_g + (\vec{r}_O-\vec{r})\times 2m\,\marc{\vec{a}_O}{N}{}\\
&= -2mga\sin\theta\,\vec{e}_z - 2ma\,\ddot{x}\cos\theta\,\vec{e}_z,
\end{aligned}$$
$$2mR^2\,\ddot{\theta} = -2mga\sin\theta - 2ma\,\ddot{x}\cos\theta,$$
\respuesta{$\displaystyle \ddot{\theta}=\frac{a}{R^2}\Big[-g\sin\theta-\ddot{x}\cos\theta\Big]$}

\textbf{Derivada del momento angular respecto a un punto acelerado}\\
El momento angular respecto a un punto $Q$ se define como
$$\marc{\vec{H}_Q}{N}{}=(\vec{r}-\vec{r}_Q)\times m(\marc{\vec{v}}{N}{}-\marc{\vec{v}_Q}{N}{})$$
Si derivamos este vector:
$$\der{\marc{\vec{H}_Q}{N}{}}{N}{}=\der{(\vec{r}-\vec{r}_Q)}{N}{}\times m(\marc{\vec{v}}{N}{}-\marc{\vec{v}_Q}{N}{})+(\vec{r}-\vec{r}_Q)\times m\,\der{(\marc{\vec{v}}{N}{}-\marc{\vec{v}_Q}{N}{})}{N}{}$$
El primer término es el producto cruz de un vector consigo mismo, luego es nulo:
$$\der{\marc{\vec{H}_Q}{N}{}}{N}{}=(\vec{r}-\vec{r}_Q)\times m(\marc{\vec{a}}{N}{}-\marc{\vec{a}_Q}{N}{})=(\vec{r}-\vec{r}_Q)\times \vec{F}-(\vec{r}-\vec{r}_Q)\times m\,\marc{\vec{a}_Q}{N}{}$$
$$\der{\marc{\vec{H}_Q}{N}{}}{N}{}=\vec{M}_Q-(\vec{r}-\vec{r}_Q)\times m\,\marc{\vec{a}_Q}{N}{}$$

\subsubsection*{Ecuaciones diferenciales del movimiento}
Eliminando $\ddot{x}$ mediante la expresión anterior:
$$\left(1+\frac{2ma^2}{(2m+M)R^2}\cos^2\theta\right)\ddot{\theta}
= \frac{a}{R^2}\left(-g\sin\theta+\frac{2ma}{2m+M}\dot{\theta}^2\sin\theta\cos\theta\right),$$
por lo que
\respuesta{$\displaystyle \ddot{\theta}
= \frac{a}{R^2}\,
  \frac{-g\sin\theta+\dfrac{2ma}{2m+M}\dot{\theta}^2\sin\theta\cos\theta}
        {1+\dfrac{2ma^2}{(2m+M)R^2}\cos^2\theta}
\qquad\text{y}\qquad
\ddot{x}=\frac{2ma}{2m+M}\Big(\dot{\theta}^2\sin\theta-\ddot{\theta}\cos\theta\Big)$}

\patotruco{Para elegir el sistema de Newton conviene mirar qué fuerzas \emph{no} actúan en cierta dirección. Con todo el conjunto (bloque + mancuerna), en la dirección horizontal no hay fuerzas externas: la primera ecuación sale gratis y sin incógnitas de reacción.}

\patoojo{Cuando tomas momento angular respecto a un punto que acelera (como $O$, que se mueve con el bloque), no basta con $\dot{\vec{H}}_O=\vec{M}_O$: aparece el término extra $-(\vec{r}-\vec{r}_O)\times m\,\vec{a}_O$. Es el que da el $-2ma\ddot{x}\cos\theta$ en la ecuación de torques.}
\end{solucion}
